% !TeX root = ../初等数学笔记.tex
% 由主文档通过 \input 载入；请编译主文档。
\setcounter{chapter}{17} % 保留原书章号。
\chapter{三角恒等变换、方程与解三角形}
\label{chap:em-triangle-solving}
三角恒等式不宜孤立记忆。和差角公式产生倍角与半角，改变加减方式又产生和积互化；这些公式与解三角形共同构成三角计算的核心。

\section{和差角公式与推导链}
\begin{theorem}[和差角公式]
对实数 $u,v$，
\begin{align}
 \cos(u\pm v)&=\cos u\cos v\mp\sin u\sin v,\label{eq:em-cos-add}\\
 \sin(u\pm v)&=\sin u\cos v\pm\cos u\sin v.\label{eq:em-sin-add}
\end{align}
\end{theorem}
\begin{proof}
单位圆上角 $u,v$ 对应点间的距离平方，一方面由坐标与勾股定理为 $2-2(\cos u\cos v+\sin u\sin v)$；另一方面旋转 $-v$ 保持距离，变为角 $u-v$ 对应点到 $(1,0)$ 的距离平方 $2-2\cos(u-v)$。相等后得到余弦差角公式，令 $v$ 换成 $-v$ 得和角式。再以 $\sin t=\cos(\pi/2-t)$ 得到正弦公式，论证没有使用后续向量数量积。
\end{proof}
在 $\cos u,\cos v,\cos(u\pm v)$ 均非零时相除，得到
\[
 \tan(u\pm v)=\frac{\tan u\pm\tan v}{1\mp\tan u\tan v}.
\]
这些分母条件是公式的组成部分，而非可忽略的附注。

\section{倍角、半角与和积互化}
令 $u=v=x$ 并使用平方关系，可得
\[
 \sin2x=2\sin x\cos x,\qquad
 \cos2x=\cos^2x-\sin^2x=1-2\sin^2x=2\cos^2x-1.
\]
于是 $\sin^2(x/2)=(1-\cos x)/2$、$\cos^2(x/2)=(1+\cos x)/2$；开方后的正负号仍由半角所在象限决定。把和角与差角式相加或相减，还得到
\[
 2\sin u\cos v=\sin(u+v)+\sin(u-v),\quad
 2\cos u\cos v=\cos(u+v)+\cos(u-v),
\]
\[
 2\sin u\sin v=\cos(u-v)-\cos(u+v).
\]
反过来令 $u=(a+b)/2,v=(a-b)/2$，便得到和差化积，例如
$\sin a+\sin b=2\sin((a+b)/2)\cos((a-b)/2)$，
$\cos a-\cos b=-2\sin((a+b)/2)\sin((a-b)/2)$。

若 $t=\tan(x/2)$ 有定义，半角式进一步给出
\[
 \sin x=\frac{2t}{1+t^2},\qquad
 \cos x=\frac{1-t^2}{1+t^2},\qquad
 \tan x=\frac{2t}{1-t^2}
\]
，最后一式还要求 $t\ne\pm1$。这种代换把三角式化为有理式，但排除了 $x=(2k+1)\pi$，解原方程时须另查这些角。

\section{三角恒等式、最值与变形策略}
恒等变形常以统一函数、降低次数或制造乘积为目标。对不同时为零的实数 $a,b$，令 $R=\sqrt{a^2+b^2}$，选择 $\varphi$ 满足 $\cos\varphi=a/R,\sin\varphi=b/R$，则
\[
 a\sin x+b\cos x=R\sin(x+\varphi).
\]
因此全实轴上的最大、最小值为 $\pm R$，且能够在相应相位取得。用 $\tan\varphi=b/a$ 单独确定角可能丢失象限，尤其在 $a=0$ 时根本无法这样除。

若 $x$ 被限制在区间，必须检查相位是否可达。例如 $\sin x+\cos x=\sqrt2\sin(x+\pi/4)$ 在 $[0,\pi/2]$ 上最大值为 $\sqrt2$，在 $x=\pi/4$ 取得；最小值为 $1$，在两个端点取得，而不是全实轴下的 $-\sqrt2$。

\section{三角方程与不等式}
对 $|c|\le1$，记 $\alpha=\arcsin c,\beta=\arccos c$，则
\[
 \sin x=c\Longleftrightarrow x=\alpha+2k\pi\ \text{或}\ x=\pi-\alpha+2k\pi,
 \quad
 \cos x=c\Longleftrightarrow x=\pm\beta+2k\pi.
\]
对任意实数 $c$，$\tan x=c$ 的解为 $x=\arctan c+k\pi$，其中 $k\in\mathbb Z$。极端值时两族表达可能重复，不产生额外解。

例如 $\sin2x=\sin x$ 化为 $\sin x(2\cos x-1)=0$，所以 $x=k\pi$ 或 $x=2k\pi\pm\pi/3$；若直接除以 $\sin x$，会漏掉整族解。三角不等式先在一个周期上确定区间，再平移：$\sin x\ge1/2$ 的解为所有 $[\pi/6+2k\pi,5\pi/6+2k\pi]$ 的并。

\section{正弦定理、余弦定理与解三角形}
设非退化三角形的边 $a,b,c$ 分别对内角 $\alpha,\beta,\gamma$，外接圆半径为 $R$。用弦长与所对圆周角关系得到\emph{正弦定理}（sine rule）
\[
 \frac a{\sin\alpha}=\frac b{\sin\beta}=\frac c{\sin\gamma}=2R.
\]
对锐角可把弦的一半放入以 $R$ 为斜边的直角三角形；钝角用补角正弦相同归约。把两边从共同顶点投影到适当坐标轴，再用勾股定理，可得\emph{余弦定理}（cosine rule）
\begin{equation}\label{eq:em-cosine-rule}
 c^2=a^2+b^2-2ab\cos\gamma,\qquad
 S=\frac12ab\sin\gamma.
\end{equation}
后一个面积式来自高为 $b\sin\gamma$；钝角时垂足在延长线上，但高仍为正。若记半周长 $s=(a+b+c)/2$，由 $16S^2=4a^2b^2-(a^2+b^2-c^2)^2$ 分解可得 Heron 公式 $S=\sqrt{s(s-a)(s-b)(s-c)}$。

三边已知先用余弦定理；两角一边已知先补第三角再用正弦定理；两边与夹角已知先求第三边。若已知 $a,b,\alpha$，则 $\sin\beta=b\sin\alpha/a$。该比值大于 $1$ 无解，等于 $1$ 时只有直角候选，小于 $1$ 时有两个互补角候选；每个候选还须满足 $\alpha+\beta<\pi$，因此可能得到零、一或两个三角形。

\section{例题、习题与学习检查}
\begin{example}
已知 $a=1,b=\sqrt3,\alpha=\pi/6$，求所有可能的三角形。
\end{example}
\begin{solve}
$\sin\beta=\sqrt3/2$，所以 $\beta=\pi/3$ 或 $2\pi/3$。两者都满足内角和限制，对应 $\gamma=\pi/2$ 或 $\pi/6$。再由 $c=\sin\gamma/\sin\alpha$ 得 $c=2$ 或 $1$，故确有两解。
\end{solve}
\begin{enumerate}
\item 求 $\sin(\pi/12)$。答：$(\sqrt6-\sqrt2)/4$。
\item 解 $\cos2x=0$。答：$x=\pi/4+k\pi/2$，$k\in\mathbb Z$。
\item 两边长为 $3,4$，夹角为 $60^\circ$，求第三边与面积。答：$\sqrt{13}$ 与 $3\sqrt3$。
\end{enumerate}

